\documentclass[../../main.tex]{subfiles}

\begin{document}

\section{Solutions for Note 3}
\begin{enumerate}
\item
If $A$ is invertible, then $A^{-1} A x = A^{-1} b$ and $x = A^{-1} b$.
Therefore, the system $Ax = b$ is always consistent.

\item
Consider the following representation of the system.
\begin{align*}
    x + 2y &= 1 \\
    2x + 4y &= k
\end{align*}
Notice that if $k = 2$, then there exist infinitely many solutions
while $k \neq 2$ gives no solutions. Thus the system never has a
unique solution.

\item
First and foremost, consider the following augmented matrix.
\[
\begin{bNiceArray}{ccc|c}
    1  & 1  & 1  & 10 \\
    3  & 4  & -1 & 2 \\
    -1 & -3 & 1  & 4
\end{bNiceArray}
\]
Using elementary row operations, the augmented matrix can be
transformed as following.
\begin{align*}
    \begin{bNiceArray}{ccc|c}
        1  & 1  & 1  & 10 \\
        3  & 4  & -1 & 2 \\
        -1 & -3 & 1  & 4
    \end{bNiceArray}
    &\rightarrow
    \begin{bNiceArray}{ccc|c}
        1  & 1  & 1  & 10 \\
        0  & 1  & -4 & -28 \\
        -1 & -3 & 1  & 4
    \end{bNiceArray}
    \rightarrow
    \begin{bNiceArray}{ccc|c}
        1 & 1  & 1  & 10 \\
        0 & 1  & -4 & -28 \\
        0 & -2 & 2  & 14
    \end{bNiceArray} \\
    &\rightarrow
    \begin{bNiceArray}{ccc|c}
        1 & 1 & 1  & 10 \\
        0 & 1 & -4 & -28 \\
        0 & 0 & -6  & -42
    \end{bNiceArray}
\end{align*}
Therefore, $-6z = -42$ and $z = 7$. Moreover, $y - 4z = -28$ and
$y = 0$. Finally, $x + y + z = 10$ and $x = 3$. Therefore, the
solution to the system is $(3, 0, 7)$.

\item
First and foremost, let $A =
\begin{bmatrix}
    1  & 1  & 1 \\
    3  & 4  & -1 \\
    -1 & -3 & 1
\end{bmatrix}$.
Using elementary row operations, $A$ can be transformed as following.
\[
\begin{bmatrix}
    1  & 1  & 1 \\
    3  & 4  & -1 \\
    -1 & -3 & 1
\end{bmatrix}
\rightarrow
\begin{bmatrix}
    1  & 1  & 1 \\
    0  & 1  & -4 \\
    -1 & -3 & 1
\end{bmatrix}
\rightarrow
\begin{bmatrix}
    1 & 1  & 1 \\
    0 & 1  & -4 \\
    0 & -2 & 2
\end{bmatrix}
\rightarrow
\begin{bmatrix}
    1 & 1 & 1 \\
    0 & 1 & -4 \\
    0 & 0 & -6
\end{bmatrix}
\]
In other words, $A$ can be represented as the matrix multiplication
below.
\begin{align*}
    \begin{bmatrix}
        1  & 1  & 1 \\
        3  & 4  & -1 \\
        -1 & -3 & 1
    \end{bmatrix}
    &=
    \begin{bmatrix}
        1 & 0 & 0 \\
        3 & 1 & 0 \\
        0 & 0 & 1
    \end{bmatrix}
    \begin{bmatrix}
        1  & 0 & 0 \\
        0  & 1 & 0 \\
        -1 & 0 & 1
    \end{bmatrix}
    \begin{bmatrix}
        1 & 0  & 0 \\
        0 & 1  & 0 \\
        0 & -2 & 1
    \end{bmatrix}
    \begin{bmatrix}
        1 & 1 & 1 \\
        0 & 1 & -4 \\
        0 & 0 & -6
    \end{bmatrix} \\
    &=
    \begin{bmatrix}
        1 & 0 & 0 \\
        3 & 1 & 0 \\
        0 & 0 & 1
    \end{bmatrix}
    \begin{bmatrix}
        1  & 0  & 0 \\
        0  & 1  & 0 \\
        -1 & -2 & 1
    \end{bmatrix}
    \begin{bmatrix}
        1 & 1 & 1 \\
        0 & 1 & -4 \\
        0 & 0 & -6
    \end{bmatrix} \\
    &=
    \begin{bmatrix}
        1  & 0  & 0 \\
        3  & 1  & 0 \\
        -1 & -2 & 1
    \end{bmatrix}
    \begin{bmatrix}
        1 & 1 & 1 \\
        0 & 1 & -4 \\
        0 & 0 & -6
    \end{bmatrix}
\end{align*}
Consider the following system for equation for $Ly = b$.
\[
\begin{bmatrix}
    1  & 0  & 0 \\
    3  & 1  & 0 \\
    -1 & -2 & 1
\end{bmatrix}
\begin{bmatrix}
    x' \\
    y' \\
    z'
\end{bmatrix}
=
\begin{bmatrix}
    10 \\
    2 \\
    4
\end{bmatrix}
\]
It is evident that $x' = 10$, $y' = 2 - 3x' = -28$, and $z' =
4 + x' + 2y' = -42$. Continuing, the system for $Ux = y$ can be
solved.
\[
\begin{bmatrix}
    1 & 1 & 1 \\
    0 & 1 & -4 \\
    0 & 0 & -6
\end{bmatrix}
\begin{bmatrix}
    x \\
    y \\
    z
\end{bmatrix}
=
\begin{bmatrix}
    10 \\
    -28 \\
    -42
\end{bmatrix}
\]
Therefore, $z = 7$, $y = 0$, and $x = 3$ holds and the solution
for the system is $(3, 0, 7)$ as shown in the previous problem.

\item
To find the inverse of $A =
\begin{bmatrix}
    1  & 1  & 1 \\
    3  & 4  & -1 \\
    -1 & -3 & 1
\end{bmatrix}$, consider the following augmentation and
transformations.
\begin{align*}
    \begin{bNiceArray}{ccc|ccc}
        1  & 1  & 1  & 1 & 0 & 0 \\
        3  & 4  & -1 & 0 & 1 & 0 \\
        -1 & -3 & 1  & 0 & 0 & 1
    \end{bNiceArray}
    &\rightarrow
    \begin{bNiceArray}{ccc|ccc}
        1  & 1  & 1  & 1  & 0 & 0 \\
        0  & 1  & -4 & -3 & 1 & 0 \\
        -1 & -3 & 1  & 0  & 0 & 1
    \end{bNiceArray} \\
    &\rightarrow
    \begin{bNiceArray}{ccc|ccc}
        1 & 1  & 1  & 1  & 0 & 0 \\
        0 & 1  & -4 & -3 & 1 & 0 \\
        0 & -2 & 2  & 1  & 0 & 1
    \end{bNiceArray} \\
    &\rightarrow
    \begin{bNiceArray}{ccc|ccc}
        1 & 1 & 1  & 1  & 0 & 0 \\
        0 & 1 & -4 & -3 & 1 & 0 \\
        0 & 0 & -6 & -5 & 2 & 1
    \end{bNiceArray} \\
    &\rightarrow
    \begin{bNiceArray}{ccc|ccc}
        1 & 1 & 1  & 1  & 0 & 0 \\
        0 & 1 & -4 & -3 & 1 & 0 \\
        0 & 0 & 1  & \frac{5}{6} & -\frac{1}{3} & -\frac{1}{6}
    \end{bNiceArray} \\
    &\rightarrow
    \begin{bNiceArray}{ccc|ccc}
        1 & 1 & 1 & 1  & 0 & 0 \\
        0 & 1 & 0 & \frac{1}{3} & -\frac{1}{3} & -\frac{2}{3} \\
        0 & 0 & 1 & \frac{5}{6} & -\frac{1}{3} & -\frac{1}{6}
    \end{bNiceArray} \\
    &\rightarrow
    \begin{bNiceArray}{ccc|ccc}
        1 & 0 & 0 & -\frac{1}{6} & \frac{2}{3}  & \frac{5}{6} \\
        0 & 1 & 0 & \frac{1}{3}  & -\frac{1}{3} & -\frac{2}{3} \\
        0 & 0 & 1 & \frac{5}{6}  & -\frac{1}{3} & -\frac{1}{6}
    \end{bNiceArray} \\
\end{align*}
Therefore, $A^{-1} =
\begin{bmatrix}
    -\frac{1}{6} & \frac{2}{3}  & \frac{5}{6} \\
    \frac{1}{3}  & -\frac{1}{3} & -\frac{2}{3} \\
    \frac{5}{6}  & -\frac{1}{3} & -\frac{1}{6}
\end{bmatrix}$ and the system could be solved by multiplying
$A^{-1}$ and $b$.
\[
\begin{bmatrix}
    -\frac{1}{6} & \frac{2}{3}  & \frac{5}{6} \\
    \frac{1}{3}  & -\frac{1}{3} & -\frac{2}{3} \\
    \frac{5}{6}  & -\frac{1}{3} & -\frac{1}{6}
\end{bmatrix}
\begin{bmatrix}
    10 \\
    2 \\
    4
\end{bmatrix}
=
\begin{bmatrix}
    3 \\
    0 \\
    7
\end{bmatrix}
\]
Thus, the solution is $(3, 0, 7)$, which is consistent with the
previous problems.
\end{enumerate}
\end{document}


